% SOURCE docs/release/2026-10-29/source/docs/writing-repair/2026-10-29/src/ad84e7f88d7e93fe.tex CHAR 0
\chapter{From an offer effect to annual contribution}
\label{rewrite:deposit}


% SOURCE docs/release/2026-10-29/source/docs/writing-repair/2026-10-29/src/ec8ce8fc42140471.tex CHAR 0
\newcommand{\DepositEstimate}{1,183.47}
\newcommand{\DepositLow}{1,128.34}
\newcommand{\DepositHigh}{1,238.60}
\newcommand{\DepositNet}{2.55}
\newcommand{\DepositNetLow}{1.03}
\newcommand{\DepositNetHigh}{4.06}


% SOURCE docs/release/2026-10-29/source/docs/writing-repair/2026-10-29/src/ad84e7f88d7e93fe.tex CHAR 136


\section{Fix the intervention before estimating its value}
The preceding accounting calculation defines the question for the first
complete study. Let $Z_i$ denote assignment of the offer to eligible household
$i$, $R_i$ its delivery, $D_i$ acceptance, and $r_i$ the paid annual rate.
These variables describe successive events, not synonyms. Let $B_i(z)$ be
the household's average eligible deposit balance over the following year under
assignment $z$, and let $C_i(z)$ be contribution over that same year. For a
population fixed before assignment, the policy targets are
\[
 \tau_B=\E[B_i(1)-B_i(0)],\qquad
 \tau_C=\E[C_i(1)-C_i(0)].
\]
A statement about the effect among acceptors generally answers neither target.
Acceptance can depend on intentions that also determine balances. Dividing an
assignment effect by an acceptance rate additionally requires an instrumental
variable argument and changes the population to the relevant compliers; the
derivation and restrictions are in Chapter~\ref{r3:instruments}.

Here the constructed design makes a narrower intervention explicit. Every
assigned offer is delivered, and assignment automatically reprices the eligible
balance for one year. Thus $R=Z$ and the price rule is $r=r_0+Z\Delta r$;
acceptance does not gate the payment. A voluntary activation campaign would
need a different expense equation. Assume no interference across households,
no lost outcomes, and unchanged funding margin and direct cost. If the
expected no-offer balance is $B_0$, the incremental expected contribution is
\begin{equation}
 \tau_C=(m-\Delta r)\tau_B-\Delta r B_0-c.
 \label{rewrite:value-map}
\end{equation}
To obtain this expression, subtract $mB_i(0)$ from
$(m-\Delta r)B_i(1)-c$, then take expectations. It charges the rate
increase on the balance that would have remained without the intervention.
The familiar accounting shortcut is justified by this particular price rule.
It would overcharge a campaign that paid only newly acquired funds and would
misdescribe a campaign whose delivery itself failed.

\section{A deliberately transparent assignment design}
The example has two equally weighted customer strata, $X=0,1$. Assignment is
independent of potential balances within each stratum, with probabilities
$e_0=0.5$ and $e_1=0.1$. Both strata have mean no-offer balance
\pounds10,000. The stipulated balance effects are $b_0=0$ and $b_1=2,400$,
with independent normal household noise of standard deviation \pounds400.
These are generated customers. In particular, the independence and normal
law are supplied by construction and have not been inferred from bank data.

Consistency, assignment exchangeability within $X$, and positive assignment
probabilities give
\[
 \tau_B=\sum_{x=0}^1 w_x
   \{\E[B\mid Z=1,X=x]-\E[B\mid Z=0,X=x]\},
 \qquad w_0=w_1=\tfrac12.
\]
The two weights describe the target population, not the observed distribution
of offered customers. The population contrast is therefore \pounds1,200.
A saturated cell-mean estimator uses the same weights. It is a useful first
implementation because every contribution to the answer can be inspected;
with two strata there is no need to learn a complex outcome surface.

A constant residual projection gives a different average when effects vary.
Its weights are proportional to $w_xe_x(1-e_x)$, giving weight
$0.09/(0.25+0.09)$ to the high-benefit stratum and a population coefficient
of about \pounds635.29. That is a valid projection target. Calling it the
equal-stratum population effect would change the financial answer from
positive \pounds3 to about negative \pounds12.53. Chapter~\ref{scisyn:evidence}
retains the separately executed repeated experiments on this target error
and on sparse support. The current single sample demonstrates the full report;
it does not add another calibration experiment.

\section{Carry uncertainty through the expense equation}
Write $n_{xa}$, $\bar B_{xa}$ and $s_{xa}$ for the count, mean and sample
standard deviation of each observed cell. For a total error allowance
$\alpha=0.05$, construct four mean intervals
\[
 I_{xa}=\bar B_{xa}\ \pm\
 t_{n_{xa}-1,1-\alpha/8}\frac{s_{xa}}{\sqrt{n_{xa}}}.
\]
Each interval has conditional noncoverage $\alpha/4$ under its normal
independent sampling law. The union bound implies simultaneous coverage
at least $1-\alpha$, without assuming independence among the four interval
events. Conditioning on the observed assignment counts is legitimate here
because assignment is independent of the outcome errors within strata.
A cell with fewer than two observations returns unsupported rather than a
zero standard error.

If $h_{xa}$ denotes a cell half-width, interval arithmetic yields
\[
 \widehat\tau_B=\tfrac12(\bar B_{01}-\bar B_{00})+
                     \tfrac12(\bar B_{11}-\bar B_{10}),\qquad
 I_B=\widehat\tau_B\ \pm\ \tfrac12\sum_{x,a}h_{xa}.
\]
This construction permits different cell variances. It is conservative and
can become too wide to guide action. The earlier generated study shows why
a narrower pooled interval needs an additional common-variance assumption.
A shared service should expose that tradeoff, not silently choose the
narrower result. The present argument does not justify the Student intervals
under heavy tails, common calendar shocks or dependent households.

With $m=0.03$, $\Delta r=0.0025$, $B_0=10,000$ and $c=5$, equation
\eqref{rewrite:value-map} is increasing in $\tau_B$. Its endpoints
therefore transform directly: $I_C=0.0275 I_B-30$. Uncertainty in these
accounting inputs would need an additional joint set or sensitivity analysis.
Here they are fixed illustrative quantities. The deterministic seed and
8,000 generated observations give the following actual calculation.
\begin{center}
\begin{tabular}{lrr}
\toprule Quantity & Estimate & Interval\\\midrule
Annual average balance effect (\pounds) & \DepositEstimate & [\DepositLow, \DepositHigh]\\
Annual contribution per household (\pounds) & \DepositNet & [\DepositNetLow, \DepositNetHigh]\\
\bottomrule
\end{tabular}
\end{center}
The four cells contain 2,039 unoffered and 1,961 offered households in the
first stratum, and 3,624 unoffered and 376 offered households in the second.
The positive contribution interval supports the action under this declared
model and expense rule. Its margin is small: an additional certain expense
above the lower endpoint removes that interval-based assurance. The result
is unsuitable as a bank forecast; its value is to make the target, support,
uncertainty and accounting inseparable in the completed report.

\section{What the implemented example actually returns}
The local callable estimator takes the three-column generated data, target
weights, four accounting inputs and $\alpha$. It returns cell counts and
uncertainty, the balance target, its interval, contribution and its interval,
and a conditional decision. The same calculation is available at the command
line from the repository root:
\begin{lstlisting}
python3 tools/run_rewrite_examples_20261029.py --report-only
\end{lstlisting}
The callable result and a separate command-line process agree exactly in
the retained execution. The generated data, code, seed and numerical results
are supplied with the study. This is a specific estimator implemented for
the example; it is not evidence that an external causal package was qualified.

The earlier MCP walkthrough establishes graph syntax, conditional
adjustment, information-date checks and refusal of a dependent sampling law.
The public study in Chapter~\ref{gap29:public} additionally sends actual
fitted estimates and covariance summaries through library, CLI and MCP.
It preserves the resulting interval and conditional or contradicted
assignment status. Neither path exposes the deposit raw-data estimator
as an MCP operation.
Chapter~\ref{scisyn:investment} gives the observed requests and the initial
timeout as well as the successful retry. These separate operations do not establish a raw-data-to-financial-report
MCP endpoint.
A future adapter must carry the complete report, including assumptions and
unsupported outcomes, before that claim becomes available.

\section{From the household offer to broader funding evidence}
Two extensions show what this transparent report leaves to be learned.
Chapter~\ref{gap29:public} estimates an assignment contrast from the retained
public experiment, compares adjusted and unadjusted estimators, examines
influence, and subtracts an explicit cost. Its earnings units and population
are different from the deposit account: the exercise tests an inferential
and reporting route, not transport of an earnings effect into bank balances.
Chapter~\ref{gap29:funding} follows institutions rather than households and
asks whether fixed effects can repair a misassigned outcome. They cannot
recover the institution's identity merely by absorbing a persistent level.
The structural funding and calendar arguments explain the additional shock
and dependence assumptions needed for that journey.

For the offer study, the next source of information is assignment support and
correct household aggregation; for the institutional panel it is a coherent
exposure and outcome history. A bank cannot substitute a successful estimate
in one of these populations for the missing design evidence in the other.
The shared report can preserve both questions without pretending they have
the same estimator or sampling law.


